% Lösungen Einheit 5 von 5 – Verteilung im Diagramm (zusammenbau v0.1)
\einheitenkopf{Einheit 5 von 5 · Verteilung im Diagramm}

\erg{22}{a) Fehlschüsse – die 3 steht als Exponent bei der Fehlschusswahrscheinlichkeit \quad b) $P(X = 2) \approx 33\,\%$ \quad c) $13$: Erwartungswert 13,5; $P(X = 13) \approx 0{,}1175$ ist größer als $P(X = 14) \approx 0{,}1141$ \quad d) Falsch: der Erwartungswert ist $40 \cdot 0{,}3 = 12$; dort liegt das Maximum, nicht bei der Hälfte der Versuche. \quad e) Abbildung 2: das Maximum muss beim Erwartungswert 3 liegen (Abbildung 1 hat es bei 5), und die Säulen müssen zusammen 100 Prozent ergeben (in Abbildung 3 ist allein die Säule bei 3 fast halb so hoch, die Summe liegt weit darüber). \quad f) $P(Y = 6) = P(X = 2) \approx 26\,\%$ \quad g) $P(X \le 3) \approx 74\,\%$ \quad h) $1 - 2 \cdot 0{,}1938 = 0{,}6124$ (wegen der Symmetrie ist $P(X \ge 8) = P(X \le 4)$) \quad i) $P(X = 7) \approx 0{,}85 - 0{,}15 - 0{,}5 = 0{,}20$: Symmetrie um 7,5 gibt $P(X = 6) = P(X = 9)$ und $P(X \ge 8) = 0{,}5$}
\erg{23}{a) $k = 12$: Erwartungswert 12; wegen der Symmetrie liegt $P(X \le 11)$ knapp unter einem Halb (0{,}4194), mit der Säule bei 12 knapp darüber (0{,}5806)}
\erg{24}{a) Fehler: das Maximum liegt beim Erwartungswert $40 \cdot 0{,}2 = 8$, nicht bei der Hälfte der Versuche. \quad b) Jede mögliche Trefferzahl hat genau eine Säule, und eine davon tritt sicher ein; die Ereignisse schließen sich aus, ihre Wahrscheinlichkeiten ergeben zusammen das sichere Ereignis. \quad c) $6$: Erwartungswert 6; $P(X = 6) \approx 0{,}1631$ ist etwas größer als $P(X = 5) \approx 0{,}1622$ \quad d) Säulen bei 0 bis 4 mit 6,25, 25, 37,5, 25 und 6,25 Prozent}

24d)

\saeulen{0/6.25,1/25,2/37.5,3/25,4/6.25}{40}{10}{$P(X = k)$ in \%}

